\section{Lesson 3: Where Is the Real Claim?}

You can now state a passage's \textbf{Claim} in plain language without dropping qualifications, fix the \textbf{Meaning} of the words doing inferential work, and inspect the \textbf{Warrant} for each premise without mistaking a famous name for a proof. Yet a sharper difficulty remains. Five confident sentences follow one another in the same declarative voice: one defines a term, the next draws a corollary, the third appeals to observation, the fourth cites an earlier chapter, and somewhere in the sequence a sentence quietly asserts something contestable about reality. Until you can put your finger on that sentence, you cannot judge what you are being asked to believe.

Finding that contestable sentence is the work of the fourth question---\textbf{Burden}---governed by a single rule: \emph{definition is not discovery}.

\subsection{Unpacking Concepts and Discovering Reality}

One distinction governs the whole craft of reading metaphysics: \emph{unpacking our concepts is not discovering something about reality}.

Both activities are necessary. Serious philosophy cannot begin unless we fix what our terms mean and trace what follows from them. The hazard is not technical language, but the unannounced \emph{transition}: the moment an argument stops explaining its words and begins telling you how the world is constituted, without flagging the shift or arguing for the new commitment.

Make one question automatic whenever you face a compressed argument:

\begin{quote}
Where does this argument stop explaining its terms and begin telling me something about reality?
\end{quote}

Our fourth question---\textbf{Burden}---turns that inquiry into a concrete target, the \textbf{burden-bearing premise}: the smallest proposition or step on which the philosophical result turns. Ask it in plain English:

\begin{quote}
What is the first proposition here whose truth is not supplied merely by vocabulary or deferred to a prior result?
\end{quote}

Working sentence by sentence, most lines fix a word, draw a corollary from what has been granted, or cite another text. Somewhere among them---stated openly, tucked inside a clause, or left implicit between two lines---sits a substantive claim about reality that the argument requires and nothing in the vocabulary supplies. That sentence is the burden. Three outcomes are possible: a single burden-bearing premise (narrowing the whole dispute to one point), several independent burdens (each named, since denying any one blocks the conclusion), or \emph{none at all} (everything follows from definitions or citations---an exact finding, not a failure of reading).

To test whether a candidate premise bears a real burden, apply the \textbf{hostile-reader test}:

\begin{quote}
Could an intelligent opponent accept all of these definitions and still deny the substantive conclusion?
\end{quote}

This opponent understands the terms, grants every distinction, follows valid logic, and still says no---stated in the strongest form his ablest defender would accept. If the answer is \emph{yes}, the argument carries genuine substantive content; find the premise he refuses. If \emph{no}, the conclusion was already contained in the definitions; say plainly that a vocabulary has been unpacked, and then ask whether the substantive claim has migrated inside the definitions themselves.

\subsection{A Worked Contrast: Potency and Change}

To measure the distance between conceptual unpacking and ontological commitment, take a standard definition from the Aristotelian and Thomistic tradition:

\begin{quote}
Potency means a capacity for an actuality not presently possessed.
\end{quote}

From this it follows immediately that:

\begin{quote}
A thing with potency has a capacity for an actuality not presently possessed.
\end{quote}

That inference is impeccable and completely empty. It reports what we have decided to mean by a word. Anyone who accepts the definition must grant the consequence, and granting it costs nothing, because nothing has been claimed about any existing thing.

Now set a second sentence beside it:

\begin{quote}
The capacity invoked by change corresponds to a real principle or feature of the changing thing, and is not merely a description of possible future states.
\end{quote}

Here we cross into a substantive claim about reality. Much of the classical tradition holds it true for serious reasons, yet it does not follow from the definition of \emph{potency}. An intelligent opponent can accept every definition word for word and still deny this sentence, holding instead that talk of capacities is a compact way of describing what a thing will or may do under conditions.

We do not settle here whether that realist claim about potency is true; our job is to mark where the argument becomes contestable. Both sentences are equally lucid English. The distance between them is not the gap between obscurity and clarity, but the gap between a claim about our words and a claim about the world.

\subsection{How Definitions Smuggle Claims and How Hidden Premises Enter}

The hardest cases are valid deductions whose definitions quietly pack substantive claims inside them. Definitions smuggle commitments in three recurring ways.

First is \textbf{the loaded definition}: a term is defined so as to presuppose the ontology at issue. Defining change as ``the actualization of a potency'' is not a neutral report of an observation; it embeds an act-and-potency analysis inside the opening word, so any act-and-potency conclusion drawn from it assumes what was at stake. The remedy is not banning the definition, but relocating the burden onto accepting that definition as true of real change.

Second is \textbf{the reified abstraction}: a phrase that begins as a way of speaking about a thing quietly becomes the name of a constituent inside it---``a capacity'' turns into ``a principle in the subject,'' or ``a way of being said'' becomes ``a mode of being.'' That step may be defensible, yet it is a substantive step taken between two sentences. Whenever an abstract noun carries heavy weight, ask whether it describes what a thing can do or names a real component inside it, and check whether the text has switched.

Third is \textbf{the stipulated bridge}: a definition connects conceptual vocabulary to reality by fiat---for example, ``by \emph{real} we mean whatever is required to make change intelligible.'' Anything reached across a stipulated bridge is reached trivially, doing by vocabulary what an argument was promised to earn.

Two companion traps belong beside these. First, distinguish a \textbf{nominal definition} (what a word signifies in speech, such as ``phoenix'' or ``first mover'') from a \textbf{real definition} (which claims to state the essence of an existing kind). Where a real definition is offered, asserting that such an essence actually exists is itself a substantive burden. Second, watch for \textbf{system-internal repetition}. Beyond crude appeals to authority (``Aquinas says \linline{P}; therefore \linline{P}''), a mature school can circle quietly inside its own terms: ``\linline{P} follows from act and potency; why accept that principle? Because that is what act and potency mean.'' Test that circle rather than dismissing it: either \linline{P} is a conceptual theorem about the vocabulary, or the ontology has migrated into the definitions and must be examined there.

Finally, compressed classical prose routinely omits steps trained readers supplied by reflex. To recover \textbf{hidden premises}, ask what the inference requires that the text leaves out, state that missing proposition in the author's idiom where possible, \textbf{mark it explicitly as your reconstruction} rather than the author's own sentence, and inspect \emph{its} warrant.

\subsection{Why This Does Not Make Metaphysics Merely Verbal}

Insisting that definition is not discovery does not make metaphysics mere word-play. Three guards keep that error away.

First, \textbf{definitions are not cheating}. Fixing a term's sense is the condition of arguing cleanly about anything; an argument that defines its terms precisely and then makes a substantive claim has done its job properly.

Second, \textbf{conceptual results can be genuinely valuable}. Showing that two accepted beliefs conflict, exposing a hidden incoherence, or proving that a distinction is needed to state a problem accurately is a real gain, though it discovers no new object in the world. From Aristotle's definition of demonstration as a syllogism productive of scientific knowledge through the cause (\emph{Posterior Analytics} I.2, 71b17--25), it follows conceptually that a valid syllogism from true premises may still fail to be a demonstration---a standard that sharpens every proof we read.

Third, \textbf{the classical tradition itself knows the difference}. Aristotle separates \emph{hypotheses}, which assert that a subject exists, from \emph{definitions}, which assert no existence (\emph{Posterior Analytics} I.2, 72a8--25), and insists we cannot grasp what a thing is without grasping that it is (\emph{Posterior Analytics} II.8, 93a16--24). Commenting on those texts in his \emph{Expositio libri Posteriorum Analyticorum}, Aquinas notes that \emph{whether} a thing is (\emph{an est}) precedes \emph{what} it is (\emph{quid est}), ``for there are no definitions of non-beings'' (I, lect.\ 2, nn.\ 4--5), and that there is no demonstration of the essence (\emph{quod quid est}) itself (II, lect.\ 7, n.\ 9). Likewise, in \emph{Summa Theologiae} I, q.\ 2, a.\ 2, ad 2, he teaches that to prove something exists we must begin from what its name signifies (\emph{quid significet nomen}), never its essence. Our six questions are a modern pedagogical aid, never attributed to Aristotle or Aquinas; yet separating nominal meaning from real existence is a discipline they practised deliberately, as our first worked passage shows.

\subsection{Worked Passage 1: Definitional in Appearance, Ontological in Fact}

Our first worked passage comes from St.\ Thomas Aquinas, \emph{Summa Theologiae} I, q.\ 2, a.\ 1 (``Whether the Existence of God Is Self-Evident?''), in the 1920 translation of the Fathers of the English Dominican Province. Here are the load-bearing sentences of the body (\emph{corpus}):

\begin{quote}
A thing can be self-evident in either of two ways: on the one hand, self-evident in itself, though not to us; on the other, self-evident in itself, and to us. A proposition is self-evident because the predicate is included in the essence of the subject, as ``Man is an animal,'' for animal is contained in the essence of man. [...] Therefore I say that this proposition, ``God exists,'' of itself is self-evident, for the predicate is the same as the subject, because God is His own existence as will be hereafter shown (Q.\ 3, Art.\ 4). Now because we do not know the essence of God, the proposition is not self-evident to us; but needs to be demonstrated by things that are more known to us, though less known in their nature---namely, by effects.
\end{quote}

Earlier in the article, Objection 2 argued that as soon as anyone understands the word ``God'' to signify ``that thing than which nothing greater can be conceived,'' God must exist in reality as well as mentally, since existing in both is greater than existing mentally alone. Note accurately that although readers recognize that formula (\emph{id quo maius cogitari non potest}) from St.\ Anselm's \emph{Proslogion}, Aquinas does \emph{not} name Anselm anywhere in \emph{ST} I, q.\ 2, a.\ 1. Here is Aquinas's Reply to Objection 2 in the same translation:

\begin{quote}
Perhaps not everyone who hears this word ``God'' understands it to signify something than which nothing greater can be thought, seeing that some have believed God to be a body. Yet, granted that everyone understands that by this word ``God'' is signified something than which nothing greater can be thought, nevertheless, it does not therefore follow that he understands that what the word signifies exists actually, but only that it exists mentally. Nor can it be argued that it actually exists, unless it be admitted that there actually exists something than which nothing greater can be thought; and this precisely is not admitted by those who hold that God does not exist.
\end{quote}

\subsubsection*{Claim}
Aquinas asserts that ``God exists'' is self-evident in itself because God is his own existence, yet is \emph{not} self-evident to us because we do not know God's essence and must demonstrate his existence from effects; and in Reply 2 he denies that unpacking ``God'' as ``something than which nothing greater can be thought'' establishes existence in reality (\emph{in rerum natura}) unless such a being is already granted to exist.

\emph{Plain-language reading.} Anyone who directly grasped God's real nature would see immediately that his nature is to exist. Because human minds do not grasp God's nature directly, ``God exists'' is not obvious to us and must be proved from his effects. Hence unpacking the definition of the word ``God'' tells us only what is in our thought, never what exists outside the mind, unless the existence in dispute has already been conceded.

\subsubsection*{Meaning}
\emph{Self-evident} (\emph{per se notum}) does not mean ``obvious''; it means the predicate belongs to the real essence of the subject. \emph{In itself} (\emph{quoad se}) names that structural relation in reality; \emph{to us} (\emph{quoad nos}) adds that our minds actually know those essences. \emph{Essence} names the real nature of a thing, not a dictionary gloss. Reply 2 keeps \emph{mental existence} (being conceived in thought) strictly apart from \emph{actual existence} in reality (\emph{in rerum natura}).

\subsubsection*{Warrant}
The account of self-evidence is \textbf{definition}. Splitting \emph{in itself} from \emph{to us} follows as a \textbf{logical consequence} once human cognitive limits are noted (illustrated in the omitted lines by common \textbf{experience} of whole and part). Crucially, ``because God is His own existence as will be hereafter shown (Q.\ 3, Art.\ 4)'' stands here strictly as a deferred \textbf{prior result}---a promissory note not redeemed on this page. That we do not know God's essence appeals to \textbf{experience} of human cognition, and Reply 2's refusal follows by \textbf{logical consequence} from separating what a word signifies from what exists in reality.

\subsubsection*{Burden}
The sentence that looks like a verbal gloss yet carries an immense ontological commitment is \emph{``God is His own existence.''} Syntactically it mirrors ``Man is an animal,'' yet it is no dictionary remark: it claims that in reality there is a being whose real essence is identical with its act of existing.

One brief pair of formulas earns its place here because it exposes the exact gap between an identity inside the signification of a name and an existential identity in reality---the slip Objection 2 commits and Aquinas blocks:
\ldisplay{(1) \lquad \logicop{Signification}(\logicop{God}) = \logicop{BeingItself}}
\ldisplay{(2) \lquad \lex x\lthin (x = \logicop{God})}
Line (1) states an identity inside the meaning of a name; line (2) asserts that something identical with God actually exists in reality. Line (1) does not entail line (2).

The outcome vindicates Aquinas's own discipline. The ontological commitment in ``God is His own existence'' is real and undischarged here---and \textbf{Aquinas blocks the hurried inference himself in Reply 2}, refusing to cross from nominal meaning to real existence and deferring both God's existence (q.\ 2, a.\ 3) and the real identity of essence and existence (q.\ 3, a.\ 4) to substantive arguments from effects. Rejecting that real identity does not yield a logical contradiction here; it commits an opponent to a \textbf{rival framework}, a distinction we will examine in Lesson 5.

\subsection{Worked Passage 2: ``Whatever Is Moved Is Moved by Another'' (Through Burden)}

We now read Aquinas's First Way in \emph{Summa Theologiae} I, q.\ 2, a.\ 3, \emph{corpus}, carrying it strictly through \textbf{Burden}; Lesson 4 will bound its \textbf{Consequence}, and Lesson 5 will assess its \textbf{Denial}. No symbolic notation appears here, because the open issue sits inside the prose itself.

Keep two Latin sentences from the Leonine edition (Rome, 1888) in view as we read: \emph{Movere enim nihil aliud est quam educere aliquid de potentia in actum} (``to move is nothing other than to bring something from potency into act'') and \emph{de potentia autem non potest aliquid reduci in actum, nisi per aliquod ens in actu} (``nothing can be brought from potency to act except by some being in act''). Note too that Aquinas's own word is \emph{viae} (``ways''), and that he develops a fuller parallel argument on ordered movers in \emph{Summa contra Gentiles} I, c.\ 13, n.\ 14. Here is the complete First Way in the 1920 English Dominican translation:

\begin{quote}
The existence of God can be proved in five ways. The first and more manifest way is the argument from motion. It is certain, and evident to our senses, that in the world some things are in motion. Now whatever is in motion is put in motion by another, for nothing can be in motion except it is in potentiality to that towards which it is in motion; whereas a thing moves inasmuch as it is in act. For motion is nothing else than the reduction of something from potentiality to actuality. But nothing can be reduced from potentiality to actuality, except by something in a state of actuality. Thus that which is actually hot, as fire, makes wood, which is potentially hot, to be actually hot, and thereby moves and changes it. Now it is not possible that the same thing should be at once in actuality and potentiality in the same respect, but only in different respects. For what is actually hot cannot simultaneously be potentially hot; but it is simultaneously potentially cold. It is therefore impossible that in the same respect and in the same way a thing should be both mover and moved, i.e.\ that it should move itself. Therefore, whatever is in motion must be put in motion by another. If that by which it is put in motion be itself put in motion, then this also must needs be put in motion by another, and that by another again. But this cannot go on to infinity, because then there would be no first mover, and, consequently, no other mover; seeing that subsequent movers move only inasmuch as they are put in motion by the first mover; as the staff moves only because it is put in motion by the hand. Therefore it is necessary to arrive at a first mover, put in motion by no other; and this everyone understands to be God.
\end{quote}

\subsubsection*{Claim}
Aquinas claims that because some things are in motion, whatever is in motion is put in motion by another, and an ordered chain of moved movers cannot stretch to infinity, it is necessary to arrive at a first mover put in motion by no other---which everyone understands to be God.

Before writing our plain-language reading, watch an instructive correction. Suppose on a first pass a reader drafts: \emph{``Everything in the world that moves physically from place to place must be pushed by something else, so there must be a first unpushed pusher.''} Checking four lines down exposes the error at once: Aquinas's own illustration is fire heating cold wood! Warming wood beside a fire does not change location or receive a mechanical push. That initial draft shrank \emph{moveri} (``being moved'') down to local locomotion, turning Aquinas's thesis into a narrower target vulnerable to an immediate appeal to inertia. Reading is iterative; catching a misreading against the author's own example is good work, so we revise openly.

\emph{Plain-language reading (revised).} Observation shows that some things undergo real change---passing from merely being able to have a feature to actually having it. Nothing can bring itself from lacking a feature to possessing it in the exact same respect; whatever changes must be actualized by something else already actual in the relevant way. Where movers act only as activated by a prior mover (like a staff moved by a hand), the series cannot stretch backward forever. Therefore there must be a first mover not moved by anything else---which people commonly call God.

\subsubsection*{Meaning}
\emph{Motion} (\emph{motus}) means any transition from potentiality to actuality, including qualitative alteration (wood warming) as well as locomotion. \emph{Potentiality} (\emph{potentia}) means a real capacity for a feature not yet possessed; \emph{actuality} (\emph{actus}) means possessing it. The qualifier \emph{in the same respect and in the same way} is vital: Aquinas allows a composite body to move itself in \emph{different} respects (an animal walks; a hot end of a rod warms its cold end), barring only simultaneous possession and lack of the identical actuality.

\subsubsection*{Warrant}
That ``some things are in motion'' appeals to \textbf{experience}. Immediately afterward, ``motion is nothing else than the reduction of something from potentiality to actuality'' enters as a loaded \textbf{definition}. That one thing cannot be simultaneously in act and potency in the same respect follows as a \textbf{logical consequence} of those definitions joined to non-contradiction.

Now inspect the hinge: ``nothing can be reduced from potentiality to actuality, except by something in a state of actuality,'' yielding ``whatever is in motion is put in motion by another.'' Two readings compete for the exact same English and Latin words, and each carries a steep price.

On the \textbf{definitional reading}, ``moved'' means passing from potency to act regarding a form, ``potency'' means lacking that form while able to have it, and ``by another'' means by something already actual in conferring it. So read, the premise states only that what comes to have a form lacked it and that nothing confers what it lacks. Its price: the premise becomes cheap and near-empty, naming no real principle inside the changing subject demanding an extrinsic cause, so the grand conclusions the tradition draws---a real composition of act and potency in every mutable thing, and an unmoved mover whose immobility is fullness of act---go unsupported.

On the \textbf{substantive reading}, the premise is a \textbf{metaphysical principle} illustrated by \textbf{experience} (fire heating wood): the potency invoked by change is a real principle inside the changing thing, not a mere description of possible future states, and it cannot actualize itself without a distinct actual mover. Its price: the conclusion now carries ontological weight, yet the premise is no longer free; it bears the philosophical burden and owes substantive defence beyond vocabulary.

Finally, excluding an infinite series of moved movers (as the staff moves only by the hand) rests on \textbf{experience} and a \textbf{metaphysical principle} compressed here as a promissory \textbf{prior result} developed in \emph{Summa contra Gentiles} I, c.\ 13, n.\ 14.

\subsubsection*{Burden}
The primary burden-bearing premise of the argument's first half is \emph{the substantive claim that the potency invoked by change is a real principle in the changing thing that cannot reduce itself to act and therefore requires a distinct actual mover}---paired, in the closing lines, with the further burden that real chains of change form an ordered series of derived movers like the staff and the hand.

Right at the center of that first burden lies an open question we deliberately leave unsettled here: \emph{is potency a real principle in the changing thing, or merely a way of describing what the thing can become under conditions?} As a description of possible future states, ``whatever becomes hot was not yet hot'' survives harmlessly but establishes no real composition demanding an extrinsic actualizer. As a real principle, the premise carries genuine weight and owes an answer to rival accounts of dispositions and physical structure. Isolating that open question cleanly prepares us for Lesson 4, where we will measure what actually follows if the burden is granted.

\subsection{Exercises}

\subsubsection*{Exercise 3.1}
Compare these two marked sentences:
\begin{enumerate}
\item (\emph{Definition}) ``An accident is a feature whose nature is to exist in a subject rather than to subsist independently on its own.''
\item (\emph{Metaphysical principle}) ``Every feature that begins to belong to a subject without flowing from that subject's own essence must be caused by an extrinsic agent.''
\end{enumerate}
Using the hostile-reader test, explain why each label is accurate and state what would change in an argument if the first were treated as a metaphysical premise or the second as a mere definition.

\subsubsection*{Exercise 3.2}
Read this short paragraph carefully:

\begin{quote}
By a \emph{contingent} entity we mean one whose non-existence involves no contradiction, and by a \emph{necessary} entity we mean one whose non-existence is impossible. Whatever is not necessary is therefore such that it could fail to exist. Moreover, since a \emph{self-explanatory} being is defined as one whose essence excludes the possibility of non-existence, no contingent entity can be self-explanatory. It follows with complete certainty that every entity is either self-explanatory or lacks an essence that guarantees its existence.
\end{quote}

What is the first proposition here whose truth is not supplied merely by vocabulary or deferred to a prior result? Justify your answer with the hostile-reader test.

\subsubsection*{Exercise 3.3}
The following short argument mimics a scholastic manual, yet suffers from a fatal defect:

\begin{quote}
Every finite being is a \emph{participated} being, where by a participated being we mean something that receives its act of existence moment by moment from an external source that does not itself receive existence from another. Since the trees, stones, and stars around us are manifestly finite beings, they are participated beings. Therefore every tree, stone, and star receives its existence moment by moment from an unreceived external source.
\end{quote}

Locate the exact defect using the hostile-reader test. Then repair the argument without smuggling the conclusion into a definition, showing that an intelligent opponent can accept every definition in your repair and still deny the substantive premise.

\subsubsection*{Exercise 3.4}
Here is an unmarked excerpt from St.\ Thomas Aquinas, \emph{Summa Theologiae} I, q.\ 3, a.\ 1, \emph{corpus} (``Whether God Is a Body?'', translated by the Fathers of the English Dominican Province, 1920):

\begin{quote}
It is absolutely true that God is not a body; and this can be shown in three ways. First, because no body is in motion unless it be put in motion, as is evident from induction. Now it has been already proved (Q.\ 2, A.\ 3), that God is the First Mover, and is Himself unmoved. Therefore it is clear that God is not a body. Secondly, because the first being must of necessity be in act, and in no way in potentiality. For although in any single thing that passes from potentiality to actuality, the potentiality is prior in time to the actuality; nevertheless, absolutely speaking, actuality is prior to potentiality; for whatever is in potentiality can be reduced into actuality only by some being in actuality. Now it has been already proved that God is the First Being. It is therefore impossible that in God there should be any potentiality. But every body is in potentiality because the continuous, as such, is divisible to infinity; it is therefore impossible that God should be a body.
\end{quote}

Locate the first fresh substantive ontological commitment on this page whose truth is not supplied by vocabulary or a backward citation, note any hidden premise the opening argument needs in this English translation, and explain why at least two earlier sentences or clauses are \emph{not} that burden.

\subsubsection*{Exercise 3.5}
Here is an unannotated passage from St.\ Thomas Aquinas, \emph{Summa Theologiae} I, q.\ 2, a.\ 3, \emph{corpus} (the Fourth Way, translated by the Fathers of the English Dominican Province, 1920). Nothing is supplied. Carry out \textbf{Claim} (with a plain-language reading), \textbf{Meaning}, \textbf{Warrant}, and \textbf{Burden} unaided:

\begin{quote}
The fourth way is taken from the gradation to be found in things. Among beings there are some more and some less good, true, noble and the like. But \emph{more} and \emph{less} are predicated of different things, according as they resemble in their different ways something which is the maximum, as a thing is said to be hotter according as it more nearly resembles that which is hottest; so that there is something which is truest, something best, something noblest and, consequently, something which is uttermost being; for those things that are greatest in truth are greatest in being, as it is written in \emph{Metaph.}\ ii. Now the maximum in any genus is the cause of all in that genus; as fire, which is the maximum heat, is the cause of all hot things. Therefore there must also be something which is to all beings the cause of their being, goodness, and every other perfection; and this we call God.
\end{quote}

\subsection{Solutions}

\subsubsection*{Solution 3.1}
Sentence 1 is a \textbf{definition}: it fixes what the Aristotelian term ``accident'' means (a non-subsisting feature such as colour or shape). No opponent can accept that meaning and deny that an accident exists in a subject; accepting it commits him to nothing about whether real things divide into substances and accidents. Treating it as a metaphysical premise wastes effort debating a verbal convention---unless an author offers it as a \emph{real} definition asserting that accidents exist in nature, which introduces a separate existential burden.

Sentence 2 is a \textbf{metaphysical principle}: it asserts a substantive causal law. An intelligent opponent can accept every definition of ``feature,'' ``subject,'' ``essence,'' and ``extrinsic agent'' and still reject it, holding instead that some features arise stochastically or without an extrinsic efficient cause. Treating Sentence 2 as a mere definition (stipulating that ``extrinsic agent'' just means whatever accounts for a new feature) turns any causal argument built on it into a tautology.

\subsubsection*{Solution 3.2}
There is \textbf{no substantive burden-bearing premise anywhere in this passage}. Every line unpacks the stipulated definitions of \emph{contingent}, \emph{necessary}, and \emph{self-explanatory}. Sentence 2 unpacks the negation of ``non-existence is impossible''; Sentence 3 draws the corollary of defining ``self-explanatory'' as excluding non-existence; Sentence 4 states a logical tautology. No opponent can accept the definitions and deny the conclusion, and accepting the whole paragraph commits nobody to asserting that anything contingent or self-explanatory exists. Manufacturing a burden here---say, crowning ``no contingent entity can be self-explanatory'' as a metaphysical discovery---mistakes conceptual bookkeeping for discovery. The paragraph keeps three terms sharp, yet discovers nothing about reality.

\subsubsection*{Solution 3.3}
The defect is \textbf{definitional smuggling} through a loaded definition and a stipulated bridge. By calling every finite being ``participated'' and defining ``participated'' to include the entire conclusion---receiving existence moment by moment from an unreceived external source---the opening sentence forces the conclusion by vocabulary alone. Under the hostile-reader test, no opponent who grants the definitions can deny the conclusion.

Here is an honest repair: ``Let a \emph{finite} being mean one limited in nature, duration, or powers. Trees, stones, and stars are finite beings. No finite being possesses existence through its own essence alone; rather, its limitation shows that its existence is conferred by an external cause. Because an essentially ordered series of causes of existence cannot stretch to infinity, every finite being depends on at least one external cause that does not receive existence from another.'' An intelligent opponent can grant the definition of ``finite'' and agree that stones and stars are finite, yet still deny that limitation entails causal dependence in existence (or contest the regress step). The substantive claim now stands openly in view.

\subsubsection*{Solution 3.4}
Two earlier clauses are \emph{not} the fresh burden here. First, ``It is absolutely true that God is not a body'' is the conclusion being argued for, not a premise. Second, ``Now it has been already proved (Q.\ 2, A.\ 3), that God is the First Mover, and is Himself unmoved'' (like ``God is the First Being'' below) is a \textbf{prior result}: a backward citation whose burden belongs to Question 2, Article 3 rather than being newly incurred on this page.

The first fresh substantive ontological commitment is the opening premise of the first argument: \emph{``no body is in motion unless it be put in motion, as is evident from induction.''} Inspecting this 1920 English translation also exposes a hidden premise: ``no body is in motion unless put in motion'' plus ``God is unmoved'' does not strictly rule out an unmoving body sitting still! To make the step valid, we supply a reconstructed premise (marked explicitly as our reconstruction): \emph{no body can move other things without itself being moved}. That substantive physical-metaphysical claim, offered on induction from experience, is the first fresh burden an opponent can contest even after granting Question 2, Article 3. In the second argument, the independent fresh burdens are the absolute priority of actuality over potentiality and the claim that every continuous body is infinitely divisible and therefore in potentiality.

\subsubsection*{Solution 3.5}
\textbf{Claim.} Aquinas claims that because beings exhibit degrees of goodness, truth, and nobility, there exists a maximum in truth, goodness, nobility, and being, which causes being, goodness, and every perfection in all other beings---and this we call God.

\emph{Plain-language reading.} Things around us are more or less good, true, and noble. Saying things have degrees of these perfections requires an actually existing maximum that they approach, as hotter things approach what is hottest; hence there is something supremely true, good, noble, and real. Because the maximum in any class causes all lesser members of that class (as fire causes heat in hot things), there must be a cause of existence, goodness, and every perfection in all beings, which people call God.

\textbf{Meaning.} \emph{More and less} concerns transcendental perfections admitting degrees across kinds, not mere size. \emph{Maximum} means an actually existing superlative reality, not a conceptual ideal. \emph{Genus} is used broadly for a domain sharing a perfection (not strictly, since being is not a genus). \emph{Cause} shifts from a standard of comparison to an efficient source.

\textbf{Warrant.} That things vary in goodness, truth, and nobility rests on \textbf{experience}. That degrees require an actually existing maximum is a \textbf{metaphysical principle} illustrated by talk of heat. Citing \emph{Metaphysics} II for ``those things that are greatest in truth are greatest in being'' supplies \textbf{provenance only} on this page (a promissory \textbf{prior result} until checked). That a maximum causes every lesser member is a second \textbf{metaphysical principle} illustrated by obsolete physics about fire. The final line is offered as a \textbf{logical consequence} plus a naming.

\textbf{Burden.} The passage rests on two independent burdens: \emph{that predicating degrees of goodness, truth, and being requires an actually existing maximum rather than a comparative scale or conceptual standard}; and \emph{that a maximum in a domain must be the efficient cause of every lesser instance in it}. Neither follows from vocabulary, from observing that one thing is better than another, or from citing \emph{Metaphysics} II.
